\section{M\"obius inversion and the exact Mellin pole terms}
\label{sec:mobius-endpoint}

The theta map in Section~\ref{sec:theta-coercivity} has an explicit
arithmetic inverse. This inverse identifies the endpoint quantities
controlled by the all-window lower bound and the quantities still
present in the codimension-one restriction.

Write \(\mu_{\mathrm M}(n)\) for the M\"obius function. For a compact
smooth function \(f\) on the physical line, define
\begin{equation}\label{eq:mobius-map}
 (\mathfrak M f)(u)
  =\sum_{n\ge1}\mu_{\mathrm M}(n)(nu)^{-1/2}f(\log(nu)),
 \qquad u>0.
\end{equation}
The sum is finite for each \(u\). If \(f\) is supported in
\((-L,L)\), then \(\mathfrak M f\) vanishes for \(u>e^L\).
With
\[
 (\Theta\psi)(x)=e^{x/2}\sum_{n\ge1}\psi(ne^x),
 \qquad
 (\operatorname{Mel}\psi)(s)=\int_0^\infty\psi(u)u^{s-1}\,du,
\]
the exact inverse and Mellin identities are
\begin{equation}\label{eq:mobius-inverse}
 \Theta\mathfrak M f=f,
 \qquad
 (\operatorname{Mel}\mathfrak M f)(s)
    =\frac{\widehat f(i(s-1/2))}{\zeta(s)}
       \quad(\Re s>1).
\end{equation}
Indeed, both compositions expand with coefficient
\(\sum_{n\mid k}\mu_{\mathrm M}(n)\), which is one for \(k=1\)
and zero otherwise. For \(\Theta\mathfrak M f\) the double sum is
finite. Conversely, for a continuous \(\psi\) decreasing faster
than every power at infinity,
\begin{equation}\label{eq:mobius-other-inverse}
 \mathfrak M\Theta\psi=\psi.
\end{equation}
For each fixed \(u>0\), its rearrangement is justified by
\(\sum_{k\ge1}d(k)|\psi(ku)|<\infty\); no estimate at \(u=0\)
is needed for this pointwise assertion. To prove the Mellin identity,
substitute \(v=nu\) in every term. For \(\sigma=\Re s>1\),
the sum of the absolute integrals is bounded by
\[
 \left(\sum_{n\ge1}n^{-\sigma}\right)
 \int_0^\infty |f(\log v)|v^{\sigma-3/2}\,dv<\infty.
\]
The remaining Dirichlet series is \(1/\zeta(s)\), and the remaining
integral is \(\widehat f(i(s-1/2))\).

\subsection{Weighted physical Plancherel on every positive shift}

For every \(\epsilon>0\), the same compact smooth \(f\) satisfies
the exact identity
\begin{equation}\label{eq:mobius-physical-plancherel}
 \int_0^\infty |\mathfrak M f(u)|^2u^{2\epsilon}\,du
 =\frac1{2\pi}\int_{\mathbb R}
 \left|\frac{\widehat f(t+i\epsilon)}
 {\zeta(1/2+\epsilon-it)}\right|^2dt<\infty.
\end{equation}
Here the zero-location theorem supplies the smoothed estimate
\begin{equation}\label{eq:mobius-smoothed}
 \mathfrak M f(u)=O_{f,\delta}(u^{-1/2-\delta})
                         \quad(u\downarrow0)
\end{equation}
for every fixed \(\delta>0\). We give a contour proof to make the
precise input explicit.

First, for fixed \(0<\delta<1/2\),
\begin{equation}\label{eq:mobius-reciprocal-growth}
 |\zeta(\sigma+it)|^{-1}
 \le C_\delta(2+|t|)^{C_\delta},
            \qquad 1/2+\delta\le\sigma\le2.
\end{equation}
The elementary continuation
\[
 \zeta(s)=\frac{s}{s-1}
           -s\int_1^\infty\{x\}x^{-s-1}\,dx
                   \qquad(\Re s>0, s\ne1)
\]
follows by integrating the integer counting function. It bounds
\(|\zeta(s)|\) by \(C(2+|t|)\) on the disks below when
\(|t|\ge4\). Take the disk centered at \(2+it\) with radius
\(R=3/2-\delta/2\). It lies in \(\Re s>1/2\), has no zero
by the zero-location theorem, and contains no pole. Choose the
analytic branch of \(\log\zeta\) whose central value is given
by the absolutely convergent Euler product and hence is bounded.
Its real part is at most \(C+\log(2+|t|)\). The
Borel--Carath\'eodory inequality on the smaller disk of radius
\(r=3/2-\delta\) gives
\(|\log\zeta(s)|\le C_\delta\log(2+|t|)\) there.
This proves \eqref{eq:mobius-reciprocal-growth} for \(|t|\ge4\).
On the remaining compact set \(1/\zeta\) is continuous, including
its removable zero at \(s=1\), which proves the bound everywhere.

On \(\Re s=2\), ordinary Mellin inversion applies to
\(\mathfrak M f\). One direct justification is that
\(x\mapsto e^{2x}\mathfrak M f(e^x)\) is smooth, vanishes for
\(x>L\), and, with all its derivatives, is \(O_f(e^x)\) as
\(x\to-\infty\), by the absolute finite sum in
\eqref{eq:mobius-map}. Fourier inversion therefore gives
\[
 \mathfrak M f(u)=\frac1{2\pi i}
 \int_{2-i\infty}^{2+i\infty}
       \frac{\widehat f(i(s-1/2))}{\zeta(s)}u^{-s}\,ds.
\]
Repeated integration by parts in the compact physical variable
makes \(\widehat f(i(s-1/2))\) decrease faster than every power
of \(|\Im s|\), uniformly on
\(1/2+\delta\le\Re s\le2\).
Equation~\eqref{eq:mobius-reciprocal-growth} thus permits shifting
the contour to \(\Re s=1/2+\delta\); the horizontal integrals
tend to zero and no pole is crossed. The new vertical integral
is absolutely convergent. Taking its absolute value proves
\eqref{eq:mobius-smoothed}.

The bound for every \(\delta>0\) makes the Mellin integral
holomorphic in \(\Re s>1/2\). Analytic continuation from
\eqref{eq:mobius-inverse} identifies it there with the same
quotient. For fixed \(\epsilon>0\), choose
\(0<\delta<\epsilon\) in \eqref{eq:mobius-smoothed}; together
with the upper support bound this gives square integrability
with weight \(u^{2\epsilon}\). Finally apply ordinary Fourier
Plancherel to
\(x\mapsto e^{(1/2+\epsilon)x}\mathfrak M f(e^x)\).
The substitution \(u=e^x\), followed by \(t=-\Im s\), gives
exactly \eqref{eq:mobius-physical-plancherel}.

The pole-order mechanism behind M\"obius mean squares is classical;
it underlies the weak Mertens criterion
\(\int_1^X(M(t)/t)^2dt=O(\log X)\), discussed with its
zero-multiplicity consequences in
\cite[preprint \S1, equations~(1.14)--(1.15), p.~3]{goldston-suriajaya2020}.
Here its exact compact smoothing is visible directly. Set
\[
 M(x)=\sum_{n\le x}\mu_{\mathrm M}(n),\qquad
 \mathfrak m(y)=e^{-y/2}M(e^y),
\]
where \(M(x)=0\) for \(0<x<1\). For every real \(Y\),
\begin{equation}\label{eq:mobius-mertens-smoothing}
 g_f(Y):=e^{-Y/2}\mathfrak M f(e^{-Y})
 =\int_{\mathbb R}\mathfrak m(Y+z)
                  \bigl(f(z)/2-f'(z)\bigr)\,dz.
\end{equation}
Indeed, Stieltjes partial summation gives
\[
 \mathfrak M f(u)
 =-u^{-1/2}\int_0^\infty M(v)v^{-3/2}
          \bigl(f'(\log(vu))-f(\log(vu))/2\bigr)\,dv.
\]
The boundary terms vanish, since the test is compactly supported
away from zero and infinity in \(v\). Substitute \(u=e^{-Y}\)
and \(v=e^{Y+z}\) to obtain
\eqref{eq:mobius-mertens-smoothing}. Consequently the left side of
\eqref{eq:mobius-physical-plancherel} is exactly
\(\int_{\mathbb R}e^{-2\epsilon Y}|g_f(Y)|^2dY\).
This identifies the smoothing and all its normalization factors;
no weak Mertens mean-square estimate is assumed here.
Minkowski's inequality gives the further exact norm comparison
\[
 \|e^{-\epsilon Y}g_f(Y)\|_2
 \le\left(\int_{\mathbb R}|f(z)/2-f'(z)|e^{\epsilon z}\,dz\right)
       \|e^{-\epsilon Y}\mathfrak m(Y)\|_2,
\]
with extended nonnegative norms allowed on the right. This follows
by translating \(Y+z\) in each norm in
\eqref{eq:mobius-mertens-smoothing}. A bound of order
\(\epsilon^{-1}\) for the squared Mertens norm is the additional
mean-square estimate at issue; the fixed-shift contour argument
above supplies no such uniform estimate.

\subsection{Local Abel masses at a zeta zero}

Let \(\gamma>0\) be a zeta ordinate, put
\(\rho=1/2-i\gamma\), and write \(m=m_\gamma\).
For any nonzero entire function \(F\), suppose its order at
\(\gamma\) is \(k<m\), and put \(r=m-k\). The meromorphic
function
\[
 Q_F(s)=\frac{F(i(s-1/2))}{\zeta(s)}
\]
has a pole of order \(r\) at \(\rho\), with leading coefficient
\begin{equation}\label{eq:mobius-pole}
 a_{\gamma,F}
 =\frac{i^k m!F^{(k)}(\gamma)}{k!\,\zeta^{(m)}(\rho)}.
\end{equation}
Choose \(\delta>0\) so that the closed disk of radius \(2\delta\)
about \(\rho\) contains no other zero or pole of \(\zeta\).
The following transform-side Abel limit is exact:
\begin{equation}\label{eq:mobius-abel}
 \lim_{\epsilon\downarrow0}
 \frac{\epsilon^{2r-1}}{2\pi}
 \int_{\gamma-\delta}^{\gamma+\delta}
 \left|\frac{F(t+i\epsilon)}
 {\zeta(1/2+\epsilon-it)}\right|^2dt
 =c_r|a_{\gamma,F}|^2,
 \qquad
 c_r=\frac{\Gamma(r-1/2)}{2\sqrt\pi\,\Gamma(r)}.
\end{equation}
To check the coefficient and normalization, write
\(Q_F(s)=(s-\rho)^{-r}H(s)\), with \(H\) analytic and
\(H(\rho)=a_{\gamma,F}\). Taylor expansion of the numerator
uses \(i(s-\rho)\), giving \eqref{eq:mobius-pole}.
In the integral, set \(t=\gamma+\epsilon v\). Its normalized
integrand becomes
\[
 \frac1{2\pi}
 \frac{|H(\rho+\epsilon(1-iv))|^2}{(1+v^2)^r}
       \mathbf1_{|v|<\delta/\epsilon}.
\]
The numerator is uniformly bounded on the chosen disk, and
\((1+v^2)^{-r}\) is integrable. Dominated convergence applies.
Finally, the substitution \(y=v^2\) gives
\(\int_{\mathbb R}(1+v^2)^{-r}dv
 =\Gamma(1/2)\Gamma(r-1/2)/\Gamma(r)\).
This proves the full limit using only the local meromorphic function.
For compact smooth \(f\),
\eqref{eq:mobius-physical-plancherel} consequently gives
\begin{equation}\label{eq:mobius-physical-lower}
 \liminf_{\epsilon\downarrow0}\epsilon^{2r-1}
 \int_0^\infty|\mathfrak M f(u)|^2u^{2\epsilon}\,du
 \ge c_r|a_{\gamma,\widehat f}|^2.
\end{equation}
For any finite collection of ordinates with the same surviving
pole order \(r\), choosing disjoint local intervals gives the sum
of their right sides. This finite-packet inequality uses positivity;
it requires no interchange of an infinite zero sum with a limit.

In particular, at a simple zero and with \(F(\gamma)\ne0\),
\begin{equation}\label{eq:mobius-simple-mass}
 \lim_{\epsilon\downarrow0}\frac{\epsilon}{2\pi}
 \int_{\gamma-\delta}^{\gamma+\delta}
 \left|\frac{F(t+i\epsilon)}
 {\zeta(1/2+\epsilon-it)}\right|^2dt
 =\frac12\left|\frac{F(\gamma)}{\zeta'(\rho)}\right|^2.
\end{equation}
The corresponding local Weil weight is \(m_\gamma|F(\gamma)|^2\).
Thus their precise conversion at a simple zero multiplies the Abel
mass by \(2|\zeta'(\rho)|^2\). At higher order,
\eqref{eq:mobius-pole} displays the derivative of \(F\) and the
derivative of \(\zeta\) that enter each surviving pole.

\subsection{The common-divisor factor in the M\"obius norm}
\label{sec:mobius-gcd}

The squared norm also has an exact arithmetic expansion that separates
its common-divisor contribution. For a compact smooth \(f\), put
\[
 A_\epsilon(f)=\int_0^\infty
          |\mathfrak M f(u)|^2u^{2\epsilon}\,du,
 \qquad
 C_\epsilon(r)=\int_{\mathbb R}e^{2\epsilon y}
       f(y+r/2)\overline{f(y-r/2)}\,dy.
\]
Here \(\epsilon\) may be complex: in the first integral the absolute
square is fixed and the weight is \(u^{2\epsilon}\). If
\(\operatorname{supp}f\subset(-L,L)\), then
\(C_\epsilon(r)=0\) for \(|r|\ge2L\). On each compact set of
\(\epsilon\)'s the functions \(C_\epsilon\) are uniformly bounded.
For \(\Re\epsilon>1/2\), direct expansion gives
\begin{equation}\label{eq:mobius-double-sum}
 A_\epsilon(f)=\sum_{n,m\ge1}
  \frac{\mu_{\mathrm M}(n)\mu_{\mathrm M}(m)}
       {(nm)^{1/2+\epsilon}}
       C_\epsilon\bigl(\log(n/m)\bigr).
\end{equation}
To justify this identity before rearranging its terms, the absolute
integral of the term indexed by \((n,m)\) is bounded by
\(B_{f,\epsilon}(nm)^{-1/2-\Re\epsilon}\).
Indeed, set \(y=\log u+(\log n+\log m)/2\) in that integral.
The two test arguments become \(y+\tfrac12\log(n/m)\) and
\(y-\tfrac12\log(n/m)\), while the weight contributes
\((nm)^{-\epsilon}e^{2\epsilon y}\).
The sum of the stated bounds converges for \(\Re\epsilon>1/2\),
which proves both termwise integration and absolute convergence of
\eqref{eq:mobius-double-sum}.

Write \(n=da\), \(m=db\), with \((a,b)=1\). The exact identity
\[
 \mu_{\mathrm M}(da)\mu_{\mathrm M}(db)
 =\mu_{\mathrm M}(a)\mu_{\mathrm M}(b)\mu_{\mathrm M}(d)^2
           \mathbf1_{(d,ab)=1}
\]
holds prime by prime. A square divisor of \(a\), \(b\), or \(d\),
or a prime common to \(d\) and \(ab\), makes both sides zero.
In the remaining case, \(a,b,d\) are pairwise coprime and squarefree,
and the formula is multiplicativity of the M\"obius function.
Absolute convergence now permits summing first over \(d\), giving
\begin{equation}\label{eq:mobius-gcd-factor}
 A_\epsilon(f)=
 \frac{\zeta(1+2\epsilon)}{\zeta(2+4\epsilon)}R_\epsilon(f),
 \qquad \Re\epsilon>1/2,
\end{equation}
where
\begin{equation}\label{eq:mobius-coprime-series}
 R_\epsilon(f)=
 \sum_{\substack{a,b\ge1\\(a,b)=1}}
 \frac{\mu_{\mathrm M}(a)\mu_{\mathrm M}(b)}
      {(ab)^{1/2+\epsilon}}
 C_\epsilon\bigl(\log(a/b)\bigr)
 \prod_{p\mid ab}(1+p^{-1-2\epsilon})^{-1}.
\end{equation}
The inner common-divisor series used here is
\[
 \sum_{(d,ab)=1}\frac{\mu_{\mathrm M}(d)^2}{d^{1+2\epsilon}}
 =\prod_{p\nmid ab}(1+p^{-1-2\epsilon})
 =\frac{\zeta(1+2\epsilon)}{\zeta(2+4\epsilon)}
       \prod_{p\mid ab}(1+p^{-1-2\epsilon})^{-1}.
\]
These products converge absolutely in the stated half-plane.

The smoothed bound~\eqref{eq:mobius-smoothed} makes
\(A_\epsilon(f)\) holomorphic in \(\Re\epsilon>0\).
For a compact subset of this half-plane choose \(\delta>0\) smaller
than its minimum real part. Near zero the absolute integrand is then
bounded by a constant times
\(u^{-1+2(\Re\epsilon-\delta)}\), and each derivative with respect
to \(\epsilon\) adds a power of \(2\log u\). These functions
are integrable uniformly on that compact subset. At the upper endpoint
the integrand has compact support. Differentiation under the integral
therefore proves the asserted holomorphy. Since the two zeta factors
are finite and nonzero in this half-plane,
\begin{equation}\label{eq:mobius-coprime-continuation}
 R_\epsilon(f):=
 \frac{\zeta(2+4\epsilon)}{\zeta(1+2\epsilon)}A_\epsilon(f),
 \qquad \Re\epsilon>0,
\end{equation}
is the holomorphic continuation of
\eqref{eq:mobius-coprime-series}. The ordinary series representation
has been established for \(\Re\epsilon>1/2\).
For real \(\epsilon>0\), the continuation is nonnegative, and is
strictly positive when \(f\ne0\): the identity
\(\Theta\mathfrak M f=f\) makes \(\mathfrak M f\ne0\),
and its smoothness then makes its weighted squared norm positive.

The pole of zeta at one gives the precise normalization
\begin{equation}\label{eq:mobius-coprime-asymptotic}
 R_\epsilon(f)=2\zeta(2)\epsilon\bigl(1+O(\epsilon)\bigr)
                      A_\epsilon(f)
              \qquad(\epsilon\downarrow0).
\end{equation}
Combining this identity with the local pole calculation
\eqref{eq:mobius-physical-lower}, for the surviving pole order \(r\),
gives
\begin{equation}\label{eq:mobius-coprime-pole}
 \liminf_{\epsilon\downarrow0}
      \epsilon^{2r-2}R_\epsilon(f)
 \ge 2\zeta(2)c_r\left|a_{\gamma,\widehat f}\right|^2.
\end{equation}
The same argument sums the right sides over every finite collection
of ordinates with that pole order. Thus the extracted common-divisor
factor accounts for one power of \(\epsilon^{-1}\); the signed
coprime expression in \eqref{eq:mobius-coprime-series}, continued by
\eqref{eq:mobius-coprime-continuation}, carries the displayed remaining
pole contributions. A uniform upper bound for this continuation near
zero is the arithmetic estimate still to be proved.

\subsection{Completion factors and the divided theta kernels}

Set
\[
 \mathcal A(s)=\frac12s(s-1)\pi^{-s/2}\Gamma(s/2),
 \qquad \xi(s)=\mathcal A(s)\zeta(s),
 \qquad \Xi(z)=\xi(1/2+iz).
\]
The functional equation makes \(\Xi\) even, and hence
\(\xi(s)=\Xi(i(s-1/2))\). Comparing the first nonzero
Taylor coefficients at \(\rho\) gives
\begin{equation}\label{eq:mobius-completion}
 \mathcal A(\rho)\zeta^{(m)}(\rho)
       =i^m\Xi^{(m)}(\gamma),
 \qquad
 a_{\gamma,F}
       =i^{-r}\mathcal A(\rho)
          \frac{m!F^{(k)}(\gamma)}{k!\Xi^{(m)}(\gamma)}.
\end{equation}
Here \(\mathcal A(\rho)\ne0\).

For \(1\le j\le m\), consider the divided transform
\(F_j(z)=\Xi(z)/(z-\gamma)^j\) constructed in
Section~\ref{sec:theta-coercivity}. It has order \(m-j\) at
\(\gamma\), so the quotient of derivatives in
\eqref{eq:mobius-completion} cancels exactly. In fact the global
meromorphic identity is
\begin{equation}\label{eq:mobius-divided}
 Q_{F_j}(s)=\frac{\mathcal A(s)}{[i(s-\rho)]^j},
 \qquad a_{\gamma,F_j}=i^{-j}\mathcal A(\rho).
\end{equation}
The inverse Fourier functions here have the two rapidly decreasing
physical tails already proved in that section. Their raw theta
lifts are explicitly
\[
 \psi_j(u)=\frac{(-i)^j u^{-\rho}}{(j-1)!}
       \int_u^\infty\!\bigl(\log(v/u)\bigr)^{j-1}
                         h(v)v^{\rho-1}\,dv,
\]
where \(h(v)=(4\pi^2v^4-6\pi v^2)e^{-\pi v^2}\).
Fubini on \(\Re s>1/2\) and
\(\int_0^1 x^{s-\rho-1}(-\log x)^{j-1}dx
 =(j-1)!/(s-\rho)^j\)
give the Mellin quotient in \eqref{eq:mobius-divided}.
The repeated first-order intertwining proved earlier gives their
theta images as the physical divided kernels. Then
equation~\eqref{eq:mobius-other-inverse} identifies those lifts with
their M\"obius inverses. In particular, the nonzero raw endpoint
coefficient persists at every allowed division order.
For this noncompact divided family the identities used are the
pointwise inverse \eqref{eq:mobius-other-inverse} and the local Abel
formula \eqref{eq:mobius-abel}. The physical Plancherel identity
\eqref{eq:mobius-physical-plancherel} has been proved on the compact
smooth core, and retains that hypothesis in its applications here.

For \(1\le j<m\), the sampled norm of the physical
divided kernel is zero, whereas its normalized local Abel mass is
\(c_j|\mathcal A(\rho)|^2>0\). At the terminal order \(j=m\),
the sampled norm is
\(m|\Xi^{(m)}(\gamma)/m!|^2\), as previously proved.
These formulas give the complete conversion on this family;
in particular, the endpoint coefficient is determined even when
the sampled values vanish at the repeated ordinate.

\subsection{What the proved coercivity estimate controls}

Write
\[
 E_L=C e^{2p_{\rm MT}L}L(1+\log L),
 \qquad p_{\rm MT}=d_{\rm MT}^{-1}>1,
\]
with \(C\) chosen so that Theorem~\ref{thm:cosine-lower-rate}
gives \(\|f\|_2^2\le e^{E_L}q(f)\) on the compact smooth
core supported in \((-L,L)\). Fourier differentiation and
Cauchy's inequality give, for every integer \(k\ge0\),
\begin{equation}\label{eq:mobius-jet-bound}
 |\widehat f^{(k)}(\gamma)|^2
 \le \frac{2L^{2k+1}}{2k+1}\|f\|_2^2
 \le \frac{2L^{2k+1}}{2k+1}e^{E_L}q(f).
\end{equation}
Indeed, the squared norm of \(x\mapsto(-ix)^ke^{-i\gamma x}\)
on \((-L,L)\) is exactly \(2L^{2k+1}/(2k+1)\).
For \(F=\widehat f\) with order \(k<m\), substitution into
\eqref{eq:mobius-abel} proves the quantitative endpoint bound
\begin{equation}\label{eq:mobius-coercive-abel}
 c_r|a_{\gamma,F}|^2
 \le c_r
 \left(\frac{m!}{k!|\zeta^{(m)}(\rho)|}\right)^2
 \frac{2L^{2k+1}}{2k+1}e^{E_L}q(f).
\end{equation}
The zero-dependent derivative factor is explicit, and the same
exponent \(E_L\) is retained.

Orthogonality to the real even trial \(v_L\) does not impose a
zero value on these evaluations. To see this by an exact test,
take a nonnegative, nonzero, even smooth \(\chi\) compactly
supported in \((-L,L)\), and put \(f(x)=\chi(x)\sin(\gamma x)\).
Then \(f\) is odd, so \(\langle f,v_L\rangle_2=0\), while
\[
 \widehat f(\gamma)
 =-i\int_{-L}^L\chi(x)\sin^2(\gamma x)\,dx\ne0.
\]
The same explicit test has mean zero. Thus restricting to this
codimension-one complement leaves the pole term in
\eqref{eq:mobius-coercive-abel} present.

For comparison with the proved prolate numerator, put
\(\lambda=e^L\). The all-window lower bound gives
\[
 \frac{q(v_L)}{b_L}
 \le C_\varepsilon\lambda^{9+\varepsilon}
       \exp\{-4\pi\lambda^2+E_L\}.
\]
Since \(E_L/\lambda^2
 =C\lambda^{2(p_{\rm MT}-1)}L(1+\log L)\to\infty\),
this bound supplies no decaying ratio. The exact inversion and
endpoint transport above retain both the larger exponent and the
displayed derivative factors. The remaining arithmetic estimate
is still a lower bound for \(b_L\) on the scale
\(\lambda^{9+\alpha}e^{-4\pi\lambda^2}\), with a positive
\(\alpha\) in the range proved possible earlier, or a direct
vanishing estimate for the arithmetic trace defects. Neither follows
from \eqref{eq:mobius-coercive-abel} alone.
